% Part:sets-functions-relations
% Chapter: size-of-sets
% Section: pairing-alt

\documentclass[../../../include/open-logic-section]{subfiles}

\begin{document}

\olfileid[pt]{sfr}{siz}{pai-alt}

\olsection{Uma Função de Emparelhamento Alternativa}

\begin{explain}
Há outras enumerações de $\Nat^2$ que facilitam descobrir
quais são suas inversas. Eis uma. Em vez de
visualizar a enumeração em uma matriz, comece com a lista de
inteiros positivos associados a espaços (inicialmente) vazios. Imagine
preencher esses espaços sucessivamente com pares $\tuple{n,m}$ da seguinte forma.
Começando com os pares que têm~$0$ na primeira posição (isto é, pares
$\tuple{0,m}$), coloque o primeiro (isto é, $\tuple{0,0}$) no primeiro espaço vazio,
depois pule um espaço vazio, coloque o segundo (isto é, $\tuple{0,2}$)
no próximo espaço vazio, pule outro, e assim por diante. O
começo (incompleto) de nossa enumeração agora se parece com isto
\[\small
\begin{array}{@{}c c c c c c c c c c c@{}}
\mathbf 1 & \mathbf 2 & \mathbf 3 & \mathbf 4 & \mathbf 5 & \mathbf 6 & \mathbf 7 & \mathbf 8 & \mathbf 9 & \mathbf{10} & \dots \\ \\
\tuple{0,0} &  & \tuple{0,1} &  & \tuple{0,2} &  & \tuple{0,3} & & \tuple{0,4} &  & \dots \\
\end{array}
\]
Repita isso com pares $\tuple{1,m}$ nos lugares que ainda permanecem
vazios, novamente pulando cada outro espaço vazio:
\[\small
\begin{array}{@{}c c c c c c c c c c c@{}}
\mathbf 1 & \mathbf 2 & \mathbf 3 & \mathbf 4 & \mathbf 5 & \mathbf 6 & \mathbf 7 & \mathbf 8 & \mathbf 9 & \mathbf{10} & \dots \\ \\
\tuple{0,0} & \tuple{1,0} & \tuple{0,1} &  & \tuple{0,2} & \tuple{1,1} & 
\tuple{0,3} & & \tuple{0,4} &  \tuple{1,2} & \dots \\
\end{array}
\]
Insira pares $\tuple{2,m}$, $\tuple{2,m}$, etc., da mesma forma. Nossa
enumeração completa assim começa assim:
\[\small
\begin{array}{@{}cc c c c c c c c c c@{}}
\mathbf 1 & \mathbf 2 & \mathbf 3 & \mathbf 4 & \mathbf 5 & \mathbf 6 & \mathbf 7 & \mathbf 8 & \mathbf 9 & \mathbf{10} & \dots \\ \\
\tuple{0,0} & \tuple{1,0} & \tuple{0,1} & \tuple{2,0}  & \tuple{0,2} & 
\tuple{1,1} & \tuple{0,3} & \tuple{3,0}  & \tuple{0,4} &  \tuple{1,2} & \dots \\
\end{array}
\]
Se numerarmos as células na matriz acima de acordo com essa
enumeração, não encontraremos uma linha zigue-zague arrumada, mas este
arranjo:
\[
\begin{array}{ c | c | c | c | c | c | c | c }
& \mathbf 0 & \mathbf 1 & \mathbf 2 & \mathbf 3 & \mathbf 4 & \mathbf 5 & \dots \\
\hline
\mathbf 0 & 1 & 3 & 5 & 7 & 9 & 11 & \dots \\
\hline
\mathbf 1 & 2 & 6 & 10 & 14 & 18 & \dots & \dots \\
\hline
\mathbf 2 & 4 & 12 & 20 & 28 & \dots & \dots & \dots \\
\hline
\mathbf 3 & 8 & 24 & 40 & \dots & \dots & \dots & \dots \\
\hline
\mathbf 4 & 16 & 48 & \dots & \dots & \dots & \dots & \dots \\
\hline
\mathbf 5 & 32 & \dots & \dots & \dots & \dots & \dots & \dots \\
\hline
\vdots & \vdots & \vdots & \vdots & \vdots & \vdots & \vdots & \ddots\\
\end{array}
\]
Podemos ver que os pares na linha~$0$ estão nas posições ímpares de
nossa enumeração, isto é, o par $\tuple{0,m}$ está na posição $2m+1$; os pares na
segunda linha, $\tuple{1,m}$, estão em posições cujo número é o
dobro de um número ímpar, especificamente, $2 \cdot (2m+1)$; os pares na
terceira linha, $\tuple{2,m}$, estão em posições cujo número é quatro vezes um
número ímpar, $4 \cdot (2m+1)$; e assim por diante. Os fatores de $(2m+1)$ para
cada linha, $1$, $2$, $4$, $8$, \dots, são exatamente as potências de~$2$:
$1= 2^0$, $2 = 2^1$, $4 = 2^2$, $8 = 2^3$, \dots\@ De fato, o
expoente relevante é sempre o primeiro membro do par em
questão. Assim, para o par $\tuple{n,m}$ o fator é $2^n$. Isso
nos dá a fórmula geral: $2^n \cdot (2m+1)$. Contudo, isso é um
mapeamento de pares para inteiros \emph{positivos}, isto é, $\tuple{0,0}$ tem
posição~$1$. Se quisermos começar na posição~$0$ devemos subtrair~$1$
do resultado. Isso nos dá:
\end{explain}

\begin{ex}
A função $h\colon \Nat^2 \to \Nat$ dada por
\[
h(n,m) = 2^n (2m+1) - 1
\]
é uma função de emparelhamento para o conjunto de pares de números naturais~$\Nat^2$.
\end{ex}

\begin{explain}
De acordo com nossa segunda enumeração de $\Nat^2$, o par
$\tuple{0,0}$ tem código $h(0,0) = 2^0(2\cdot 0+1) - 1 = 0$;
$\tuple{1,2}$ tem código $2^{1} \cdot (2 \cdot 2 + 1) - 1 = 2
\cdot 5 - 1 = 9$; $\tuple{2,6}$ tem código $2^{2} \cdot (2
\cdot 6 + 1) - 1 = 51$.
\end{explain}

Às vezes basta codificar pares de números naturais~$\Nat^2$
sem exigir que a codificação seja sobrejetora. Tais codificações têm
inversas que são apenas funções parciais.

\begin{ex}
A função $j\colon \Nat^2 \to \Nat^+$ dada por
\[
j(n,m) = 2^n3^m
\]
é uma função !!{injective} $\Nat^2 \to \Nat$.
\end{ex}

\end{document}
